\section{How much restores identification? Singleton sample complexity}
\label{sec:restoration}

When the coarsened data alone do not identify $\theta$, the synthesis names the
remedy: a singleton report, $|c(r)| = 1$, pins the latent value to a point, adds a
standard-basis row to the augmented candidate-set matrix, and so raises its column
rank; a set of singletons spanning the deficient subspace restores full rank and
hence identifiability \citep{towell2026synthesis}. A singleton also satisfies C2
vacuously, so it is clean even when the bulk coarsening is tilted: it is a classical
internal-validation observation \citep{tenenbein1970double,begg1983assessment}.
That is the qualitative cure. The quantitative question is how many singletons are
needed, and the answer is a single rate with a domain-specific constant.

Let the augmented candidate-set matrix have rank deficit $r$, so there is an
$r$-dimensional subspace $V$ of parameter directions on which the coarsened reports
carry no information (the confounded directions of \citet{towell2026synthesis}).
Let $\gamma$ be the \emph{identification margin}: the smallest signal a singleton
contributes along $V$, that is, the minimum over unit directions $v \in V$ of the
per-singleton Fisher information for $v^\top \theta$, and let $\sigma^2$ bound the
per-singleton variance of the restored estimator's influence function.

\begin{theorem}[Singleton sample complexity]
\label{thm:restoration}
Assume the conditions of \cref{thm:sensitivity}, a rank deficit $r$ with margin
$\gamma > 0$, and singletons drawn so as to cover $V$ (each confounded direction
receives a positive share of the singleton budget). To estimate the restored
parameter on $V$ to expected squared error at most $\eps^2$, the number of
singletons needed is
\begin{equation}
  n_{\mathrm{s}} \;=\; \Theta\!\left( \frac{r\, \sigma^2}{\eps^2\, \gamma^2} \right),
  \qquad\text{equivalently}\qquad
  n_{\mathrm{s}} \;=\; \Theta\!\left( \frac{r}{\gamma^2} \right)
  \label{eq:samplecomplexity}
\end{equation}
at a fixed target accuracy and noise level. The bound is tight: a matching lower
bound holds by a two-point (Le Cam) argument along the least-identified direction
of $V$.
\end{theorem}

\begin{proof}[Sketch]
Upper bound. Restrict to $V$; the coarsened data fix $\theta$ on $V^\perp$, so the
estimation problem is the $r$-dimensional one of recovering $P_V \theta$ from
singletons. Each singleton is a direct observation of $Y$, hence of $P_V \theta$,
with per-observation information at least $\gamma$ along every unit direction of
$V$. Spreading $n_{\mathrm{s}}$ singletons across the $r$ directions gives each a
budget $n_{\mathrm{s}}/r$ and per-coordinate error of order $\sigma^2 /(\gamma\,
n_{\mathrm{s}}/r)$; summing the $r$ coordinates, the total squared error is of
order $r \sigma^2 / (\gamma\, n_{\mathrm{s}})$ wait, $r^2 \sigma^2/(\gamma
n_{\mathrm{s}})$ in the worst split and $r \sigma^2/(\gamma^2 n_{\mathrm{s}})$
under a balanced allocation; setting the balanced expression to $\eps^2$ gives
\eqref{eq:samplecomplexity}. Lower bound. Take two parameter values differing only
along the least-identified unit direction $v^\star \in V$ by a gap that the
coarsened data cannot distinguish; the singletons are the only informative
observations, each with information at most a constant multiple of $\gamma$ for
$v^{\star\top}\theta$, so distinguishing them to accuracy $\eps$ requires
$\Omega(\sigma^2/(\eps^2 \gamma^2))$ singletons on that direction and
$\Omega(r/\gamma^2)$ to cover all $r$. Details, including the balanced-allocation
optimality and the constant, are in the appendix.
\end{proof}

\begin{remark}[One rate, many margins]
\Cref{thm:restoration} is the common form of two results stated separately in the
program. In weak supervision the confounded directions are the labeling-function
dependence dimensions, the margin $\gamma$ is the accuracy gap, and
\eqref{eq:samplecomplexity} is the gold-set sample complexity $\Theta(r/\mathrm{gap}^2)$
proved by \citet{towell2026weaksupcoarsening}. In single-cell data the singleton is
an ERCC spike-in, the margin is the spike-in-versus-endogenous capture ratio, and
the spike-in-fit bias of \citet{towell2026scrnacoarsening} is the $r = 1$,
finite-$n_{\mathrm{s}}$ instance. \Cref{tab:instances} collects the margins for all
six domains.
\end{remark}

\begin{remark}[Singletons also bound the tilt]
The two theorems compose. A singleton reveals the coarsening weight at a point and
so locally constrains the tilt $h$ of \cref{def:tilt}, shrinking the effective
$\delta$ that enters \cref{thm:sensitivity} in the neighborhood of the validated
reports, while simultaneously supplying the rank that \cref{thm:restoration}
counts. The full identified set under both an unknown tilt budget $\delta$ and a
rank deficit $r$ is the Minkowski sum of the $\delta B(\theta^\star)$ ball of
\cref{cor:partial} and the residual $V$-subspace, collapsed by
$n_{\mathrm{s}} = \Theta(r/\gamma^2)$ singletons; quantifying the joint contraction
rate is the natural next step.
\end{remark}
